\subsection{DIMENSION AND CODIMENSION OF A SUBSPACE OF A VECTOR SPACE}

PROPOSITION 4. Every subspace F of a vector space E is a direct factor of E and

\[
\dim \mathbf {F} + \dim (\mathbf {E} / \mathbf {F}) = \dim \mathbf {E}.\tag{3}
\]

As the quotient vector space E/F is a free module, we know (§ 1, no. 11, Proposition 21) that F is a direct factor of E; relation (3) is then a special case of formula (1) (no. 2).

COROLLARY 1. If E, F, G are vector spaces over a field K, every exact sequence of linear mappings \(0 \rightarrow E \rightarrow F \rightarrow G \rightarrow 0\) splits.

This is another way of expressing Proposition 4 (§ 1, no. 9).

COROLLARY 2. Let \((\mathbf{E}_i)_{0\leqslant i\leqslant n}\) be a finite family of vector spaces over a field K. If there exists an exact sequence of linear mappings

\[
0 \longrightarrow \mathrm{E} _ {0} \stackrel {u _ {0}} {\longrightarrow} \mathrm{E} _ {1} \stackrel {u _ {1}} {\longrightarrow} \mathrm{E} _ {2} \longrightarrow \dots \longrightarrow \mathrm{E} _ {n - 1} \stackrel {u _ {n - 1}} {\longrightarrow} \mathrm{E} _ {n} \stackrel {u _ {n}} {\longrightarrow} 0\tag{4}
\]

the relation

\[
\sum_ {2 k + 1 \leqslant n} \dim \mathrm{E} _ {2 k + 1} = \sum_ {2 k \leqslant n} \dim \mathrm{E} _ {2 k}\tag{5}
\]

holds, or, if all the spaces are finite dimensional,

\[
\sum_ {i = 1} ^ {n} (- 1) ^ {i} \dim \mathrm{E} _ {i} = 0.\tag{6}
\]

Let \(\mathbf{I}_k = \operatorname{Im} u_k = \operatorname{Ker} u_{k+1}\) for \(0 \leqslant k \leqslant n-1\) ; \(\mathbf{I}_{k+1}\) is therefore isomorphic to \(\mathbf{E}_{k+1}/\mathbf{I}_k\) , hence (formula (3)) \(\dim \mathbf{I}_k + \dim \mathbf{I}_{k+1} = \dim \mathbf{E}_{k+1}\) for \(0 \leqslant k \leqslant n-2\) and moreover \(\dim \mathbf{I}_0 = \dim \mathbf{E}_0\) and \(\mathbf{I}_{n-1} = \mathbf{E}_n\) , hence \(\dim \mathbf{I}_{n-1} = \dim \mathbf{E}_n\) . Replacing \(\dim \mathbf{E}_i\) by its expression as a function of the

\[
\sum_ {k = 0} ^ {n - 1}
\]

dim \(I_{k}\) in the two sides of (5), we obtain on each side \(\sum_{k=0}^{n-1}\dim I_{k}\) , whence the corollary.

COROLLARY 3. If M and N are two subspaces of a vector space E, then

\[
\dim (\mathbf {M} + \mathbf {N}) + \dim (\mathbf {M} \cap \mathbf {N}) = \dim \mathbf {M} + \dim \mathbf {N}.\tag{7}
\]

It suffices to apply Corollary 2 to the exact sequence

\[
0 \rightarrow \mathbf {M} \cap \mathbf {N} \rightarrow \mathbf {M} \oplus \mathbf {N} \rightarrow \mathbf {M} + \mathbf {N} \rightarrow 0
\]

(§ 1, no. 8, Proposition 10) taking account of the fact that

\[
\dim (\mathbf {M} \oplus \mathbf {N}) = \dim \mathbf {M} + \dim \mathbf {N}
\]

(no. 2, Proposition 2).

COROLLARY 4. For every subspace \(\mathbf{F}\) of a vector space \(\mathbf{E}\) , \(\dim \mathbf{F} \leqslant \dim \mathbf{E}\) ; if \(\mathbf{E}\) is finite dimensional, the relation \(\dim \mathbf{F} = \dim \mathbf{E}\) is equivalent to \(\mathbf{F} = \mathbf{E}\) .

The first assertion is obvious from (3); further, if dim E is finite, the relation dim F = dim E implies dim(E/F) = 0 by (3) and a vector space of dimension 0 reduces to 0.

COROLLARY 5. If a vector space \(\mathbf{E}\) is the sum of a family \((\mathbf{F}_{\iota})\) of vector subspaces, then

\[
\dim \mathrm{E} \leqslant \sum_ {\iota} \dim \mathrm{F} _ {\iota}.\tag{8}
\]

If further dim E is finite, the two sides of (8) are equal if and only if E is the direct sum of the family \((\mathbf{F}_{\iota})\) .

The inequality (8) follows from (3) and the fact that E is isomorphic to a quotient of \(\bigoplus_{\iota} F_{\iota}\) (§ 1, no. 7, formula (28)). The second assertion is a particular case of § 1, no. 10, Corollary 5 to Proposition 16, for the equality of the two sides of (8) implies that \(\dim F_{\iota} = 0\) except for a finite number of indices.

DEFINITION 2. Given a vector space E, the codimension (with respect to E) of a subspace F of E, denoted by \(\operatorname{codim}_{E} F\) , or simply codim F, is the dimension of E/F (equal to that of any supplement of F in E).

Relation (3) may then be written

\[
\dim \mathrm{F} + \operatorname{codim} \mathrm{F} = \dim \mathrm{E}.\tag{9}
\]

PROPOSITION 5. Let \(\mathbf{F}\) , \(\mathbf{F}'\) be two subspaces of a vector space \(\mathbf{E}\) , such that \(\mathbf{F} \subset \mathbf{F}'\) . Then \(\operatorname{codim}_{\mathbf{E}} \mathbf{F}' \leqslant \operatorname{codim}_{\mathbf{E}} \mathbf{F} \leqslant \dim \mathbf{E}\) . If \(\operatorname{codim}_{\mathbf{E}} \mathbf{F}\) is finite, the relation \(\operatorname{codim}_{\mathbf{E}} \mathbf{F}' = \operatorname{codim}_{\mathbf{E}} \mathbf{F}\) implies \(\mathbf{F} = \mathbf{F}'\) .

The inequality \(\operatorname{codim}_{\mathbf{E}} \mathbf{F} \leqslant \dim \mathbf{E}\) is obvious from (9) and if \(\dim \mathbf{E}\) is finite the relation \(\operatorname{codim}_{E} F = \dim E\) implies \(\dim F = 0\) and hence \(F = \{0\}\) . The remainder of the proposition follows from this, for

\[
\operatorname{codim} _ {\mathbf {E}} \mathbf {F} ^ {\prime} = \operatorname{codim} _ {\mathbf {E} / \mathbf {F}} (\mathbf {F} ^ {\prime} / \mathbf {F}),
\]

since \(\mathbf{E} / \mathbf{F}'\) is canonically isomorphic to \((\mathbf{E} / \mathbf{F}) / (\mathbf{F}' / \mathbf{F})\) (I, § 4, no. 7, Theorem 4).

PROPOSITION 6. If M and N are two subspaces of a vector space E, then

\[
\operatorname{codim} (\mathbf {M} + \mathbf {N}) + \operatorname{codim} (\mathbf {M} \cap \mathbf {N}) = \operatorname{codim} \mathbf {M} + \operatorname{codim} \mathbf {N}.\tag{10}
\]

It suffices to apply Corollary 2 of Proposition 4 to the exact sequence

\[
0 \rightarrow \mathrm{E} / (\mathrm{M} \cap \mathrm{N}) \rightarrow (\mathrm{E} / \mathrm{M}) \oplus (\mathrm{E} / \mathrm{N}) \rightarrow \mathrm{E} / (\mathrm{M} + \mathrm{N}) \rightarrow 0
\]

(§ 1, no. 7, Proposition 10) and use no. 2, Proposition 2.

Note that if \(\mathbf{E}\) is finite dimensional, (10) is a consequence of (7) and (9).

PROPOSITION 7. If \((\mathbf{F}_i)\) is a finite family of subspaces of a vector space \(\mathbf{E}\) , then \(\operatorname{codim}\left(\bigcap_i \mathbf{F}_i\right) \leqslant \sum_i \operatorname{codim} \mathbf{F}_i\) .

If \(\mathbf{F} = \bigcap_{i}\mathbf{F}_{i}\) , \(\mathbf{E} / \mathbf{F}\) is isomorphic to a subspace of the direct sum of the \(\mathbf{E} / \mathbf{F}_i\) (§ 1, no. 7, formula (27)).

Vector subspaces of dimension 1 (resp. dimension 2) of a vector space E are often called lines passing through 0 (resp. planes passing through 0) (or simply lines (resp. planes) if no confusion arises (cf.~§9, no. 3)), by analogy with the language of Classical Geometry; a subspace of E is called a hyperplane passing through 0 (or simply a hyperplane) if it is of codimension 1. Hyperplanes can also be defined as the maximal elements of the set S of vector subspaces of E distinct from E, ordered by inclusion. There is a one-to-one correspondence between the subspaces of E containing a subspace H and the subspaces of E/H (I, §4, no. 7, Theorem 4); if E is of dimension ≥1, S is non-empty and to say that H is maximal in S means that E/H contains no subspace distinct from \{0\} and E/H, which implies that E/H is generated by any of its elements ≠0, in other words it is of dimension 1.

In a vector space of finite dimension \(n \geqslant 1\) , the hyperplanes are the subspaces of dimension n - 1 by formula (3).

PROPOSITION 8. In a vector space E over a field K, every vector subspace F is the intersection of the hyperplanes which contain it.

It suffices to show that for all \(x \notin F\) there exists a hyperplane H containing F and not containing x. By hypothesis \(F \cap Kx = \{0\}\) and hence the sum M of F and Kx is direct. Let N be supplementary to M in E; E is then the direct sum of \(H = F + N\) and Kx and H is therefore a hyperplane with the desired property.

Remark. Most of the properties proved in this no. for subspaces of a vector space do not hold for submodules of a free module whose dimension (§ 2, no. 7, Remark 2) is defined. *For example, an ideal of a commutative ring does not necessarily admit a basis, for there are integral domains A in which certain ideals are not principal (VII, § 1, no. 1) and any two elements of such a ring are linearly dependent (§ 1, no. 11, Remark 1).* A submodule of a free A-module E may be free, distinct from E and have the same dimension as E, as is shown by principal ideals in an integral domain A; the same example proves moreover that a free submodule of a free A-module does not necessarily admit a supplement.

